\section{Additional results and diagnostic partitions}
\label{app:additional}

\begin{table}[h]
\centering
\caption{Ordered post-simulation evidence-gate partition.}
\label{tab:gatepartition}
\small
\begin{tabularx}{\linewidth}{@{}p{0.34\linewidth}p{0.16\linewidth}X@{}}
\toprule
First failing/succeeding condition & Count & Interpretation \\
\midrule
No strict feasible retained state & 2 & Violation-ranked fallback; cannot support strict-core evidence. \\
Classical ambiguity below .08 & 7 & Strong classical methods give little reason to focus on disagreement. \\
Concentration below frozen threshold & 14 & QAOA $p=2$ does not meet both absolute and relative mass thresholds. \\
All evidence conditions pass & 9 & Post-simulation evidence-positive subset; still not a deployment decision. \\
\bottomrule
\end{tabularx}
\end{table}


\paragraph{Depth and matched draws.}
Across instances, mean optimum masses are .022278 (uniform), .249182 ($p=1$), and .152190 ($p=2$). The $p=2$ value exceeds uniform on 30/32; $p=1$ exceeds $p=2$ on 28/32. Since the parameter searches are unequal and selected by expected energy before optimum mass, this is a warning against interpreting depth as the sole causal factor.

Uniform and $p=2$ near-optimal coverage means are .919271 and .963542. Their paired mean difference .044271 has 95\% interval [$-.020833,.111979$]. We therefore report the direction but make no positive coverage claim. The first-hit means of 36.50 and 12.9375 are sensitive to the frozen draw stream and miss code; they complement, rather than replace, the exact probability-mass result.

\paragraph{Feasibility decomposition.}
Thirty instances have at least one strict feasible retained state; two use fallback states. All 32 cluster-repair incumbents have zero explicit disallowed-mode violations, while 27 satisfy the stricter conjunction of review capacity and empirical service/capacity tolerances. Any abstract phrase ``zero final hard violations'' should be read as zero mode-domain violations only.

\paragraph{Gate interpretation.}
The ordered partition in \cref{tab:gatepartition} prevents double counting. Its 9 positive instances meet all frozen predicates after QAOA. A real invocation router would have to replace $p_2^\star$ with pre-execution features and be evaluated on new held-out instances. Until then, ``gate rate'' is an evidence summary, not a rate of saved QPU calls.
